\section{\ldeen{Serien und Standalone}{Series and Standalone}}
\label{sec:series-standalone}

\ldeen{Die Klasse \texttt{osglecture.cls} unterstützt zwei Arbeitsformen: 
Serie und Standalone. 
Die Form die wir bisher betrachtet haben, war die Serie. 
In einer Serie liefert OLLM den Buildauftrag, kennt die 
Nachbarlehreinheiten und verwaltet Referenz- und
Integrationszustände. 

Im Standalone-Betrieb beschreibt das Dokument seine
Ausgabe selbst und wird ohne Projektmanifest gesetzt. Die fachlichen Befehle
bleiben weitgehend gleich; unterschiedlich ist vor allem, woher Kontext und
Defaults stammen.}{The same class supports two workflows. In a series, OLLM
provides the build request, knows neighbouring units, and manages reference and
integration state. In standalone operation, the document describes its own
output and is typeset without a project manifest. The subject-facing commands
remain largely the same; the main difference is where context and defaults
come from.}

\subsection{Integration}
\ldeen{Die Dokumente der Lehreinheiten einer Serie können einzeln verwendet werden.
Dies ist in der Regel bei Präsentationen sinnvoll.
Sie können jedoch auch zu einem gesamten Dokument integriert werden.
So könnte jede Lehreinheit ein Kapitel eines Skripts bilden, die dann zu einem Gesamtskript
integriert werden. Es kann aber z.\,B. auch eine Handoutsammlung angelegt
werden.
}{}
\ldeen{Ein Gesamtskript entsteht in \osglecture\ in der Regel nicht durch 
\texttt{\string\include} der TeX-Quellen, sondern durch Zusammenfügen bereits
übersetzter Dokumente.}{}
\ldeen{Dafür muss es eine eigene Integrationseinheit
mit einem Rahmendokument geben. 
Diese Einheit ist durch den Buchstaben "`i"' beim Rollenelement gekennzeichnet.
DIe laufende Nummer des Verzeichnisses ist egal, weil es bei der Zählung
nicht berüksichtig wird. Es empfiehlt sich aber, hierfür entweder die kleinste
oder die größte Nummer zu wählen, so dass die zu integrierenden Einheiten 
aufeinander folgen können.
}{}

\begin{verbatim}
course/
  ollmconfig.toml
  Include/
    projectconfig.tex
    documentmetadata.tex
  010-introduction/main.tex
  020-processes/main.tex
  900-i-collection/main.tex
\end{verbatim}

\ldeen{Der normale Workflow der Integration ist inkrementell: 
Zunächst bauen Sie beim Schreiben die zu integrierenden Dokumente und erst zuletzt 
das Rahmendokument.
Es empfiehlt sich, vor dem des des Rahmendokuments mit \texttt{ollm check} zu überprüfen
ob alle zu integrierenden Dokumente erstellt und alle Referenzen aufgelöst sind.
Falls nicht, müssen Sie ggf.\ weitere Übersetzungsläufe anstossen.
Auch hier kann Sie \texttt{ollm} unterstützen: 
Mit Hilfe der Option \texttt{--resolve} werden Abhängigkeiten bis zu einer
vorgegeben Tiefe erkannt und ausgelöst.}{}

\begin{verbatim}
cd 020-processes
ollm script --language=de
ollm script --language=de --scope=series --resolve
ollm script --language=de --scope=collection --resolve
ollm check --scope=series
\end{verbatim}

\subsubsection{\ldeen{Rahmendokumentunit}{Wrapper Document}}
\ldeen{ Die Klasse erkennt diese Rolle am Verzeichnisnamen und lädt den
Integrationsdienst automatisch. Das Dokument nennt die einzubindenden
logischen Units in der gewünschten Reihenfolge}{A complete script is not
created by applying \texttt{\string\include} to the TeX sources. Instead, the
wrapper document is a separate unit with role \texttt{i}. The class recognizes
this role from the directory name and loads the integration service
automatically. The document names the logical units to include in the desired
order}:

\begin{verbatim}
\documentclass{osglecture}

\begin{document}
\includeunit{introduction}
\includeunit{processes}
\includeunit{synchronization}
\end{document}
\end{verbatim}

\ldeen{\texttt{\string\includeunit} wählt automatisch dieselbe Sprache und
denselben Dokumenttyp wie das Rahmendokument. Es importiert also bei einem
deutschen Skript genau die deutschen Skriptprojektionen. OLLM baut die
Collection mit \texttt{--scope=collection}; existieren keine oder mehrere
passende Integrationsunits für den Dokumenttyp, ist die Auswahl nicht
eindeutig und der Build schlägt fehl.}{\texttt{\string\includeunit}
automatically selects the same language and document type as the wrapper. Thus
a German script imports exactly the German script projections. OLLM builds the
collection with \texttt{--scope=collection}; if there is no matching
integration unit for the document type, or more than one, the selection is
ambiguous and the build fails.}

\ldeen{Jede einzubindende Projektion muss zuvor mindestens einmal gebaut sein.
Der Integrationslauf protokolliert die tatsächlich verwendete Generation als
Abhängigkeit. Nach Änderungen genügt deshalb normalerweise ein erneuter Lauf
mit \texttt{--resolve}. Die Reihenfolge der
\texttt{\string\includeunit}-Befehle ist die Reihenfolge im Ergebnis; sie wird
nicht aus den Verzeichnisnummern erraten.}{Every projection to be included must
have been built at least once. The integration run records the generation
actually used as a dependency. After changes, another run with
\texttt{--resolve} is therefore normally sufficient. The order of the
\texttt{\string\includeunit} commands is the order in the result; it is not
inferred from directory numbers.}

\subsubsection{\ldeen{Tags, Links und Fallback}{Tags, Links, and Fallback}}

\ldeen{Sind Rahmendokument und importierte Units als Tagged PDF gebaut,
übernimmt \texttt{tagpax} Seiten, Strukturbaum, Ziele, Links,
Inhaltsverzeichniseinträge und Lesezeichen. Dadurch bleiben
Barrierefreiheitsstruktur und Navigation über Unit-Grenzen erhalten. Dafür
muss das Target \texttt{DocumentMetadata} erlauben beziehungsweise verlangen,
und \texttt{documentmetadata.tex} muss Tagging aktivieren.}{If the wrapper and
imported units are built as Tagged PDF, \texttt{tagpax} carries across pages,
the structure tree, destinations, links, table-of-contents entries, and
bookmarks. This preserves accessibility structure and navigation across unit
boundaries. The target must permit or require \texttt{DocumentMetadata}, and
\texttt{documentmetadata.tex} must enable tagging.}

\begin{verbatim}
\DocumentMetadata{
  lang=\OsgLectureRequestedLanguage,
  tagging=on
}
\end{verbatim}

\ldeen{Für ungetaggte Projektionen verwendet die Integration einen Fallback
über \texttt{pdfpages} und \texttt{newpax}. Links und Gliederung bleiben dabei
weitgehend funktionsfähig, die importierten Seiten werden jedoch als Artefakt
eingebunden und besitzen keine erschlossene Lesereihenfolge oder
Überschriftenstruktur. Gemischte getaggte und ungetaggte Units zwingen für den
betroffenen Dokumenttyp den ungetaggten Pfad und erzeugen eine deutliche
Warnung. Eine vollständig getaggte Serie ist daher nicht nur qualitativ besser,
sondern auch vorhersehbarer.}{For untagged projections, integration uses a
fallback based on \texttt{pdfpages} and \texttt{newpax}. Links and document
navigation largely remain functional, but imported pages are included as an
artifact and have no accessible reading order or heading structure. Mixing
tagged and untagged units forces the untagged path for the affected document
type and produces a prominent warning. A consistently tagged series is
therefore both higher quality and more predictable.}

\subsection{\ldeen{Standalone}{Standalone}}

\ldeen{Ein Standalone-Dokument besitzt weder Buildauftrag noch
Serienmanifest. Es muss dies mit der Klassenoption \texttt{standalone}
ausdrücklich erklären. Dokumenttyp und Sprache, die in einer Serie von OLLM
kommen, werden nun mit \texttt{doctype} und \texttt{lang} gewählt. Mit
\texttt{profile} kann das Dokumentprofil überschrieben werden}{A standalone
document has neither a build request nor a series manifest. It must declare
this explicitly with the \texttt{standalone} class option. The document type
and language that OLLM supplies in a series are now selected with
\texttt{doctype} and \texttt{lang}. The document profile can be overridden
with \texttt{profile}}:

\begin{verbatim}
\documentclass[
  standalone,
  doctype=script,
  profile=scrbook,
  lang=de,
  class-options={11pt,oneside}
]{osglecture}
\end{verbatim}

\ldeen{\texttt{doctype} ist ohne Angabe \texttt{slides}; für die eingebauten
Dokumenttypen existieren Defaultprofile. Explizite Werte machen ein
Standalone-Dokument dennoch verständlicher und reproduzierbarer.
\texttt{class-options} reicht Optionen direkt an die vom Profil gewählte
Basisklasse weiter. Eine gleichzeitig vorhandene OLLM-Auftragsdatei und die
Option \texttt{standalone} sind ein Fehler, weil sonst zwei Quellen denselben
Buildkontext bestimmen würden.}{\texttt{doctype} defaults to \texttt{slides},
and built-in document types have default profiles. Explicit values nevertheless
make a standalone document clearer and more reproducible.
\texttt{class-options} passes options directly to the base class selected by
the profile. Having both an OLLM build request and the \texttt{standalone}
option is an error because two sources would otherwise determine the same build
context.}

\ldeen{Ohne Manifest wird \texttt{Include/projectconfig.tex} nicht automatisch
gefunden. Gemeinsam genutzte Pakete, Sprachkonfiguration, Titel und eigene
Makros müssen daher im Dokument geladen oder ausdrücklich aus einer Datei
eingelesen werden. Ein eigenständiges mehrsprachiges Skript kann beispielsweise
so beginnen}{Without a manifest,
\texttt{Include/projectconfig.tex} is not discovered automatically. Shared
packages, language configuration, title, and custom commands must therefore be
loaded in the document or explicitly read from a file. A standalone
multilingual script might begin as follows}:

\begin{verbatim}
\documentclass[standalone,doctype=script,profile=scrbook,lang=de]
  {osglecture}
\usepackage[
  languages={de,en},
  targetlang=de,
  map={de=ngerman,en=british},
  load babel
]{langselect}
\title{\ldeen{Betriebssysteme}{Operating Systems}}
\end{verbatim}

\ldeen{Profile wie \texttt{ltx-talk} verlangen
\texttt{\string\DocumentMetadata} bereits vor
\texttt{\string\documentclass}. In einem Standalone-Dokument gibt es keine
von OLLM vorgeschaltete \texttt{documentmetadata.tex}; der Aufruf muss daher
wirklich am Dateianfang stehen}{Profiles such as \texttt{ltx-talk} require
\texttt{\string\DocumentMetadata} before
\texttt{\string\documentclass}. A standalone document has no
\texttt{documentmetadata.tex} inserted by OLLM, so the call must genuinely be
placed at the start of the file}:

\begin{verbatim}
\DocumentMetadata{lang=en,tagging=on}
\documentclass[
  standalone,doctype=slides,profile=ltx-talk,lang=en
]{osglecture}
\end{verbatim}

\ldeen{Sie können ein solches Dokument direkt mit Lua\LaTeX\ beziehungsweise
\texttt{latexmk} bauen. Der kompatible Komfortaufruf
\texttt{ollm standalone main.tex} verwendet ebenfalls die Klassenoptionen;
Target-, Sprach- und Matrixoptionen von OLLM sind dabei nicht zulässig. Ohne
Serienzustand gibt es keine externen Unit-Referenzen, Counter-Continuation oder
\texttt{\string\includeunit}-Auflösung. Lokale Referenzen und alle
modusabhängigen Darstellungsbefehle funktionieren unverändert.}{Such a
document can be built directly with Lua\LaTeX\ or \texttt{latexmk}. The
compatible convenience command \texttt{ollm standalone main.tex} also uses the
class options; OLLM target, language, and matrix options are not permitted in
this mode. Without series state there are no external unit references, counter
continuation, or \texttt{\string\includeunit} resolution. Local references and
all mode-aware rendering commands continue to work unchanged.}
